\documentclass[11pt]{article}
\usepackage[T1]{fontenc}
\usepackage{amsmath,amssymb,booktabs}
\usepackage[margin=1in]{geometry}
\usepackage[hidelinks]{hyperref}
\usepackage{tabularx}
\title{Proof Dossier: Ragasa2026A Spectral Power Equals Parent Power}
\author{Project Koios Formal}
\date{}
\begin{document}
\maketitle
\section{Identity and status}
\begin{tabularx}{\textwidth}{@{}l>{\raggedright\arraybackslash}X@{}}
\toprule
Statement ID & \texttt{BG-RAGASA2026A-THM-006} \\
Status & \texttt{proved} \\
Lean declaration & \path{ProjectKoios.BandGauge.Papers.Ragasa2026A.spectralSplittingField_power_eq_parent} \\
Module & \path{ProjectKoios/BandGauge/Papers/Ragasa2026A/RankTwoFields.lean} \\
\bottomrule
\end{tabularx}

\section{Exact statement}
For the exact finite diagonalization witness,
\[
 \sum_i\frac{\alpha_i}{2E_{\mathrm{ref}}^2}
   \sum_n d_{i,n}^2
 =
 \sum_i\frac{\alpha_i}{2E_{\mathrm{ref}}^2}
   \lVert D(q_i)\rVert_F^2.
\]

\section{Source and proof}
This theorem \texttt{formalizes} the common spectral/operator power identification used
in manuscript lines 603--610 and Supplementary Section S1.1 at immutable
\path{ksdft2effmass} commit
\begin{center}\small\texttt{99a46cb59d0c6f27ada7d709f194996ef9619d9f}.\end{center}
The Frobenius entry sum of the real diagonal matrix is $d_{i,0}^2+d_{i,1}^2$.
BG-BLOCH-THM-006 transfers that power through the unitary diagonalization.

\section{Boundary and trusted base}
No numerical eigenvalue or matrix is supplied. The identity is symmetric in the two
spectral labels and requires no ordering. Lean and pinned Mathlib report
\texttt{propext}, \texttt{Classical.choice}, and \texttt{Quot.sound}. No admission or
scientific axiom is used. AI-assisted and kernel-checked; no independent expert review
is recorded. Status: \texttt{proved}; working draft.
\end{document}
